% Source: tex/chapters/ch06/06_ソフトウェアエンジニアリング運用提供.tex
\subsubsection{問題解決}

問題解決は、ソフトウェアやシステムのインシデントと問題の原因を識別し、分析し、事業への混乱を最小化する活動である。繰り返し発生する本番問題には、ソフトウェア、インフラ、設定、データベース、ネットワーク、人間の手順など、複数の原因が絡むことがある。そのため、問題解決には、ソフトウェア技術者、運用担当技術者、セキュリティ担当、データベース担当、製品責任者などの多職種チームが必要になる場合がある。

問題解決では、監視、ログ、プロファイル、実行トレース、設定履歴を使って、何が起きたか、いつ起きたか、どの変更と関係するかを調べる。問題の根本原因がわかれば、再発防止策を実装する。
